%% ====================================================================
%% Pipelines
%% ====================================================================

\section[Pipelines]{Running pipelines with Snakemake}
\begin{frame}{Why a pipeline instead of one big job script?}
\small
\begin{cmdtablex}{lLL}
                  & \textbf{One job script}                & \textbf{A Snakemake pipeline}             \\
    \midrule
    Jobs          & One big job: samples one after another & One job per step and sample, many at once \\
    A step fails  & Start everything again                 & Fix it: only failed and missing steps run \\
    Resources     & Booked for the biggest step, all along & Each step books what it needs             \\
    Keeping track & You check which samples are done       & Snakemake checks the files                \\
\end{cmdtablex}

\bigskip
The lab's pipelines are Snakemake pipelines: you run them, you do not have to write them
\end{frame}

\begin{frame}{A pipeline runs every step, for every sample}
\diagram{readqc-dag}

\bigskip
\begin{itemize}
    \item Each step's output is the next step's input: Snakemake works out the order
    \item Every step runs as its own Slurm job, many at the same time
    \item Run it again and only missing or failed steps run
    \item The lab's pipelines are Snakemake pipelines, laid out like today's example:
          \cmd{HPC_intro/readqc}
\end{itemize}
\end{frame}

\begin{frame}{Inside a pipeline directory}
\small
\DTsetlength{0.2em}{1em}{0.2em}{0.4pt}{1.6pt}
\renewcommand*\DTstyle{\ttfamily}
\renewcommand*\DTstylecomment{\rmfamily\footnotesize\color{darkgray}}
\dirtree{%
.1 readqc/.
.2 README.md\DTcomment{how to install and run it}.
.2 snakemake\_*\_conda\_spec.txt\DTcomment{to create the \texttt{snakemake} env}.
.2 \alert{config/}\DTcomment{\alert{yours: what to run on, and how}}.
.3 \alert{config.yaml}\DTcomment{the settings, e.g.\ \texttt{min\_length}}.
.3 \alert{sample\_table.csv}\DTcomment{your samples: you write this}.
.2 workflow/\DTcomment{the pipeline itself: leave it as it is}.
.3 Snakefile\DTcomment{the start: loads the rules}.
.3 rules/\DTcomment{the steps, e.g.\ \texttt{qc.smk}}.
.3 envs/\DTcomment{the tools, with their exact versions}.
.3 profiles/\DTcomment{how to run: \texttt{hpc\_slurm/} or \texttt{local/}}.
.3 schemas/\DTcomment{checks your config and sample sheet}.
.2 results/, logs/\DTcomment{appear when you run it}.
}

\bigskip
The lab's pipelines have the same layout: \alert{\texttt{config/}} is yours,
\texttt{workflow/} is the pipeline
\end{frame}

\begin{frame}[fragile]{A step is a \emph{rule}: from input files to output files}
\begin{lstlisting}[style=footnotesize]
rule seqkit_stats:
    input:
        rules.fastp.output.fwd,
        rules.fastp.output.rev,
    output:
        "results/{sample}/02_stats/{sample}_stats.tsv",
    conda:
        "../envs/qc.yaml"
    log:
        "logs/seqkit_stats/{sample}.log",
    shell:
        "seqkit stats -a -T {input} > {output} 2> {log}"
\end{lstlisting}
\hint{From \cmd{readqc/workflow/rules/qc.smk}}
\end{frame}

\begin{frame}[fragile]{The sample sheet: which reads belong to which sample}
You write it yourself: one line per sample, the columns separated by commas
\begin{lstlisting}[style=scriptsize,language={}]
sample_name,fastq_r1,fastq_r2
sample_A,../exercise_data/sample_A/sample_A_R1.fastq.gz,../exercise_data/sample_A/sample_A_R2.fastq.gz
\end{lstlisting}
\small
\begin{cmdtablex}{lL}
    \textbf{Column}     & \textbf{What goes in it}                       \\
    \midrule
    \cmd{sample_name}   & The name: it becomes \cmd{results/<sample>/}   \\
    \cmd{fastq_r1}      & Where the R1 reads are                         \\
    \cmd{fastq_r2}      & Where the R2 reads are                         \\
\end{cmdtablex}

\bigskip
Paths start from \cmd{readqc}, so \cmd{../exercise_data} is one directory up.
The first line, with the column names, must stay exactly as it is
\end{frame}

\begin{frame}[fragile]{Running a pipeline on the HPC}
In \cmd{readqc}, on the login node (fine: Snakemake only submits jobs)
\begin{lstlisting}[style=footnotesize]
conda activate snakemake       # create it once, see README

# 1. sample sheet: which reads belong to which sample
nano config/sample_table.csv

# 2. dry run: what would run? Nothing runs yet
snakemake -j unlimited -n \
    --workflow-profile workflow/profiles/hpc_slurm/

# 3. run it: Snakemake submits every step to Slurm
snakemake -j unlimited \
    --workflow-profile workflow/profiles/hpc_slurm/
\end{lstlisting}
\hint{Start it in \cmd{screen}: Snakemake keeps going until the last job is done.
    On your laptop: \cmd{--workflow-profile workflow/profiles/local/}}
\hint{Not in Slurm account \cmd{dla_mm}? Set \cmd{slurm_account} in the profile}
\end{frame}

\begin{frame}{Exercises: a Snakemake pipeline}
\begin{exlist}{5_pipelines.md}
    \ex{5.1}{Set up the pipeline}{1}
    \ex{5.2}{Dry run}{1}
    \ex{5.3}{Run it}{1}
    \ex{5.4}{Run it again}{2}
    \ex{5.5}{Change a setting}{3}
\end{exlist}
\end{frame}

\begin{frame}{What you learned today}
\begin{itemize}
    \item Move around, work with files, \cmd{find}, \cmd{grep}, loops and scripts
    \item Save your work with Git and GitHub
    \item Install tools in Conda environments
    \item Log in to the HPC, run jobs with \cmd{srun} and \cmd{sbatch}, and run a pipeline
\end{itemize}

\bigskip
Those 200 FASTQ files? Check them with \cmd{md5sum}, keep your script in Git,
install the assembler with Conda, and run it with \cmd{sbatch}
\end{frame}
